\uuid{hSL9}
\titre{Théorème de Bézout}
\auteur{bodin ; croizat ; sacre}
\organisation{exo7}
\licence{CC BY-NC-SA 4.0}
\datecreate{2022-03-08}
\niveau{L1}
\module{Algèbre}
\chapitre{Arithmétique dans Z}
\sousChapitre{Nombres premiers, nombres premiers entre eux}
\difficulte{2}
\theme{Théorème de Bézout}
\exo7qcmid{hilisit.arithmetique.021}

\contenu{
\begin{question}
\qclassification{103.04}
    Pour les entiers $a,b$ suivants, les $u,v$ donnés sont-ils des coefficients de Bézout, c'est-à-dire tels que $au+bv=\pgcd(a,b)$ ?
    \begin{answers}
        \bad{$a=7$, $b=11$, $u=2$, $v=-3$}
        \bad{$a=20$, $b=55$, $u=6$, $v=-2$}
        \good{$a=28$, $b=12$, $u=1$, $v=-2$}
        \good{$a=36$, $b=15$, $u=-2$, $v=5$}
    \end{answers}
    \begin{explanations}
    Pour $a=7$, $b=11$, $u=2$, $v=-3$, on a $\pgcd(a,b)=1$ et $au+bv= -19 \neq 1$. \\
    Pour $a=20$, $b=55$, $u=6$, $v=-2$, on a $\pgcd(a,b)=5$ et $au+bv= 10 \neq5$. \\
    Pour $a=28$, $b=12$, $u=1$, $v=-2$, on a $\pgcd(a,b)=4$ et $au+bv=4$. \\
    Pour $a=36$, $b=15$, $u=-2$, $v=5$, on a $\pgcd(a,b)=3$ et $au+bv=3$. \\
%   On note $d=\pgcd(a,b)$. Les coefficients de Bézout vérifient $au+bv=d$. \\
%    $a=7$, $b=11$, $d=1$, $u=-3$, $v=2$ \\
%    $a=20$, $b=55$, $d=5$, $u=3$, $v=-1$ \\
%    $a=28$, $b=12$, $d=4$, $u=1$, $v=-2$ \\
%    $a=36$, $b=15$, $d=3$, $u=2$, $v=-5$ \\`
    \end{explanations}
\end{question}
}
